\begin{trapbox}
\begin{description}[style=nextline, leftmargin=2.8em, labelsep=0.5em, itemsep=0.35em]
    \item[\textbf{\textcolor{CTrap}{易错点一（忽略内角和）}}] 对三个角不满足 $A+B+C=\pi$ 的数组，直接代入恒等式进行计算。
    \item[\textbf{\textcolor{CTrap}{正确理解（三角约束）：}}] 公式仅在 $A,B,C$ 为三角形内角时成立，使用前必须验证和为 $\pi$。
    \item[\textbf{\textcolor{CTrap}{错因分析（随意推广）：}}] 学生常把该式当作普适的代数恒等式，忘记背后依赖 $A+B+C=\pi$ 的推导前提。
\end{description}

\medskip
\begin{description}[style=nextline, leftmargin=2.8em, labelsep=0.5em, itemsep=0.35em]
    \item[\textbf{\textcolor{CTrap}{易错点二（符号误判）}}] 计算 $\cos A\cos B\cos C$ 时，忽略钝角余弦为负，导致结果偏大。
    \item[\textbf{\textcolor{CTrap}{正确理解（余弦正负）：}}] 乘积在锐角三角形为正，直角三角形为零，钝角三角形为负，直接影响 $\sin^2$ 和的取值范围。
    \item[\textbf{\textcolor{CTrap}{错因分析（忽视符号）：}}] 思维定势认为所有三角函数值均为正，未根据角度范围判断余弦符号。
\end{description}

\medskip
\begin{description}[style=nextline, leftmargin=2.8em, labelsep=0.5em, itemsep=0.35em]
    \item[\textbf{\textcolor{CTrap}{易错点三（平方混淆）}}] 把 $\sin^2 A+\sin^2 B+\sin^2 C$ 与 $(\sin A+\sin B+\sin C)^2$ 混为一谈，用和平方公式展开而忽略交叉项。
    \item[\textbf{\textcolor{CTrap}{正确理解（恒等关系）：}}] 平方和与和平方没有直接简化关系，必须使用本恒等式或降次、积化和差另行处理。
    \item[\textbf{\textcolor{CTrap}{错因分析（记忆混乱）：}}] 对三角恒等式中“平方和”类公式记忆不牢，混淆了 $(a+b+c)^2$ 的代数展开与本题结构。
\end{description}
\end{trapbox}
